\documentclass{article}
\usepackage{amsmath,amssymb}
\usepackage{geometry}
\geometry{a4paper, margin=1in}

\begin{document}

\section*{Week 8. Lecture 19.}

\section{Function $\text{Ln } z$}
\textbf{Def 15.1:} The natural logarithm of a complex number $z$ is a complex number $w$ such that $e^w = z$. Notation: $w = \text{Ln } z$.

The logarithm of $r \in \mathbb{R}$ denote as usual $\ln r$.

Let us develop a formula for the computation of the complex logarithm. Write: $z = r(\cos \varphi + i\sin \varphi), w = u + iv$.

By Def 15.1: $e^{u+iv} = r(\cos \varphi + i\sin \varphi) \implies e^u (\cos v + i\sin v) = r(\cos \varphi + i\sin \varphi)$.

By comparing the complex numbers, we conclude:
\begin{align*}
    &e^u = r \implies u = \ln r \\
    &\begin{cases}
    \cos v = \cos \varphi \\
    \sin v = \sin \varphi
    \end{cases} \implies v = \varphi + 2\pi k, k \in \mathbb{Z}
\end{align*}
$\implies \text{Ln } z = w = u + iv = \ln r + i(\varphi + 2\pi k)$ for $k \in \mathbb{Z}$.

Recall: $r = |z|, \varphi + 2\pi k = \arg z$, thus (*) $\text{Ln } z = \ln |z| + i \arg z$.

$\implies \text{Ln } z$ is defined for any $z \neq 0$ and it is multi-valued function (for different values of $k$ the function attains different values).

\textbf{Ex 17.1:} Compute: 1) $\text{Ln } i$ 2) $\text{Ln } (-4)$

Sol: 1) $\text{Ln } i = \ln |i| + i(\arg i + 2\pi k) = \ln 1 + i(\frac{\pi}{2} + 2\pi k) = i(\frac{\pi}{2} + 2\pi k)$.

2) $\text{Ln } (-4) = \ln |-4| + i(\arg(-4) + 2\pi k) = \ln 4 + i(\pi + 2\pi k)$.

\textbf{Ex 17.2:} Find an analytic function $f(z)$ such that $\text{Re } f(z) = \arctan \frac{y}{x}$ and $f(-i) = \frac{3\pi}{2} + 2i$.

Sol: In Ex 11.3 we have seen that $f(z) = \arctan \frac{y}{x} - \frac{i}{2} \ln(x^2 + y^2) + ic$.

$\implies$ the only question is to determine the value of the constant $c$.

Let $z = x + iy$. Recall: $\arg z - \theta$ s.t. $\cos \theta = \frac{x}{|z|}, \sin \theta = \frac{y}{|z|} \implies \frac{\sin \theta}{\cos \theta} = \tan \theta = \frac{y}{x} \implies \theta = \arctan \frac{y}{x} \implies \arg z = \arctan \frac{y}{x}$ (up to $2\pi$), $|z|^2 = x^2 + y^2$.
$\implies f(z) = \arg z - \frac{i}{2} \ln |z|^2 + ic = -i(\ln|z| + i \arg z) + ic$, or $f(z) = -i \text{Ln } z + c_1$, where $c_1 \in \mathbb{C}$ is a constant, let us compute it: $\frac{3\pi}{2} + 2i = f(-i) = -i\text{Ln }(-i) + c_1 = -i(0 + i\frac{3\pi}{2}) + c_1 = \frac{3\pi}{2} + c_1 \implies c_1 = 2i$ and $f(z) = i\text{Ln } z + 2i$.

\textbf{Prop 15.1:} For any $z_1, z_2 \in \mathbb{C}, n \in \mathbb{N}$:
\begin{enumerate}
    \item $\text{Ln}(z_1 z_2) = \text{Ln } z_1 + \text{Ln } z_2$
    \item $\text{Ln}(z_1/z_2) = \text{Ln } z_1 - \text{Ln } z_2$
    \item $\text{Ln } z^n = n \text{Ln } z$
    \item $\text{Ln} \sqrt[n]{z} = \frac{1}{n} \text{Ln } z$
\end{enumerate}

Pf: By formula (*), Prop 1.2, and Prop 1.3, we get:
$\text{Ln}(z_1 z_2) = \ln|z_1 z_2| + i\arg(z_1 z_2) = \ln(|z_1||z_2|) + i(\arg z_1 + \arg z_2) = (\ln|z_1| + i\arg z_1) + (\ln|z_2| + i\arg z_2) = \text{Ln } z_1 + \text{Ln } z_2$.

The rest in the same way. (FINISH!)

Rmk: If $z_1 = r_1 e^{i\varphi_1}, z_2 = r_2 e^{i\varphi_2}$, then 1) can be written as $\text{Ln}(z_1 z_2) = \ln r_1 + \ln r_2 + i\varphi$, where $\varphi = \varphi_1 + \varphi_2 \pmod{2\pi}$.

Q-n: CHECK - 2), 3), 4)!

Now let us study $\text{Ln } z$ in the same spirit as we studied the complex root.

$w = \text{Ln } z$: For any $z$ correspond infinite number of values $w$. We choose a single-valued branch by fixing $k_0$: $w_{k_0} = \ln r + i(\varphi + 2\pi k_0)$.

Let $z$ move on some closed curve that does not encircle the origin, $z = 0$. After completing a full circle, $z$ returns to the same point with the same values of $r$ and $\varphi \implies$ the point $w_{k_0}$ does not change.

If $z$ moves (anticlockwise) on some closed curve that encircles the origin, then when $z$ returns to the starting point the modulus $r$ is the same, but the argument of $z$ will increase by $2\pi$ and $\text{Ln } z$ will attain a new value: $w_{new} = \ln r + i((\varphi + 2\pi k_0) + 2\pi) = \ln r + i(\varphi + 2\pi(k_0 + 1)) = w_{k_0 + 1}$.

So, after completing a full circle around the origin, the branch $w_{k_0}$ "jumps" to a new branch $w_{k_0 + 1}$. If two full circles are completed: $\text{Ln } z$ will attain another (new) value: $w_{k_0 + 2}$, since the modulus is the same, but the argument will increase by $4\pi$. In the same way, by completing finite number of full circles (anticlockwise) around the origin it is possible to pass from $w_{k_0}$ branch of $\text{Ln } z$ to any given in advance branch. For example, after 7 full circles around the origin $w_{k_0}$ will pass to $w_{k_0 + 7}$ (if in anticlockwise direction) or to $w_{k_0 - 7}$ (if clockwise direction).

The point $z = 0$ is the branch point of the function $\text{Ln } z$.

Conclusion: In any domain that does not contain closed curves

\end{document}